\documentclass[a4paper,11pt]{article}
\usepackage[utf8]{inputenc}
\usepackage[T1]{fontenc}
\usepackage[french]{babel}
\usepackage{geometry}
\geometry{left=2cm,right=2cm,top=1.5cm,bottom=2cm}
\usepackage{amsmath,amssymb}
\usepackage{enumitem}
\usepackage{array}
\usepackage{booktabs}
\usepackage{xcolor}
\usepackage{tcolorbox}
\usepackage{fancyhdr}
\usepackage{tikz}
\usepackage{lastpage}

% Couleurs
\definecolor{bleupro}{HTML}{0969da}
\definecolor{orange}{HTML}{d97706}
\definecolor{vertpro}{HTML}{2f855a}
\definecolor{rouge}{HTML}{c53030}
\definecolor{gris}{HTML}{718096}

% En-tête
\pagestyle{fancy}
\fancyhf{}
\lhead{\small \textcolor{rouge}{CORRECTION} — CCF Physique-Chimie — Terminale ICCER}
\rhead{\small Fluides \& Courant alternatif}
\cfoot{\thepage/\pageref{LastPage}}
\renewcommand{\headrulewidth}{0.5pt}

% Boîtes
\tcbuselibrary{skins,breakable}

\newtcolorbox{corrbox}{
  colback=green!3,
  colframe=vertpro,
  title={\textbf{Correction}},
  fonttitle=\bfseries,
  breakable
}

\newtcolorbox{contextebox}{
  colback=blue!3,
  colframe=bleupro,
  leftrule=4pt,
  rightrule=0pt,
  toprule=0pt,
  bottomrule=0pt,
  boxsep=4pt,
  left=8pt
}

\newcommand{\pts}[1]{\hfill\textcolor{gris}{\small\textbf{#1}}}
\newcommand{\comp}[1]{\colorbox{blue!10}{\scriptsize\textbf{#1}}}

\setlength{\parindent}{0pt}
\setlength{\parskip}{6pt}

\begin{document}

\begin{center}
{\LARGE\bfseries\color{rouge} CORRECTION — CCF Physique-Chimie}\\[8pt]
{\large Terminale Bac Pro ICCER — Groupement 1}\\[8pt]
{\Large\bfseries Débit de fluide \& Puissance en régime sinusoïdal}\\[12pt]
\textcolor{rouge}{\rule{12cm}{1pt}}
\end{center}

\vspace{10pt}

% ================================================================
\section*{Exercice 1 — TP numérique : dimensionner un réseau de chauffage \pts{5 points}}
% ================================================================

\begin{contextebox}
\textbf{Données :} $Q_v = 1{,}8$ m³/h \quad $d_1 = 40$ mm \quad $d_2 = 25$ mm\\
Vitesse recommandée : entre 0,5 et 1,5 m/s
\end{contextebox}

\vspace{6pt}

% --- Q1 ---
\textbf{Q1} \comp{APP} — Relevé des valeurs de la simulation. \pts{1 pt}

\begin{corrbox}
Avec $Q_v = 1{,}8$ m³/h et $d_1 = 40$ mm, la simulation affiche :

\begin{center}
\renewcommand{\arraystretch}{1.5}
\begin{tabular}{|>{\centering\arraybackslash}m{4cm}|>{\centering\arraybackslash}m{4cm}|}
\hline
\textbf{Grandeur} & \textbf{Valeur} \\
\hline
Section $S_1$ & \textcolor{rouge}{\textbf{12,57 cm²}} \\
\hline
Vitesse $v_1$ & \textcolor{rouge}{\textbf{0,40 m/s}} \\
\hline
Norme respectée ? & \textcolor{rouge}{\textbf{Non (0,40 < 0,5 m/s)}} \\
\hline
\end{tabular}
\renewcommand{\arraystretch}{1}
\end{center}

\textit{Barème : 0,25 pt par ligne correcte + 0,25 pt pour la conclusion sur la norme.}
\end{corrbox}

% --- Q2 ---
\textbf{Q2} \comp{REA} — Vérification par le calcul. \pts{1 pt}

\begin{corrbox}
\textbf{Étape 1 — Conversion du débit :}
\[
Q_v = 1{,}8 \text{ m}^3\text{/h} = \frac{1{,}8}{3\,600} = 5{,}00 \times 10^{-4} \text{ m}^3\text{/s}
\]

\textbf{Étape 2 — Calcul de la section :}
\[
d_1 = 40 \text{ mm} = 0{,}040 \text{ m}
\]
\[
S_1 = \pi \times \left(\frac{d_1}{2}\right)^2 = \pi \times (0{,}020)^2 = \pi \times 4 \times 10^{-4} = 1{,}257 \times 10^{-3} \text{ m}^2
\]

\textbf{Étape 3 — Calcul de la vitesse :}
\[
v_1 = \frac{Q_v}{S_1} = \frac{5{,}00 \times 10^{-4}}{1{,}257 \times 10^{-3}} = \boxed{0{,}40 \text{ m/s}}
\]

La valeur calculée confirme celle de la simulation \checkmark

\smallskip
\textit{Barème : 0,25 pt conversion + 0,25 pt section + 0,5 pt vitesse.}
\end{corrbox}

% --- Q3 ---
\textbf{Q3} \comp{ANA} — Relevé des branches ($d_2 = 25$ mm). \pts{1 pt}

\begin{corrbox}
Avec $d_2 = 25$ mm, la simulation affiche :

$v_2 = \boxed{0{,}51 \text{ m/s}}$

\textbf{Vérification :}
\begin{itemize}[nosep]
  \item $Q_{v2} = Q_v / 2 = 5{,}00 \times 10^{-4} / 2 = 2{,}50 \times 10^{-4}$ m³/s
  \item $S_2 = \pi \times (0{,}0125)^2 = 4{,}91 \times 10^{-4}$ m²
  \item $v_2 = 2{,}50 \times 10^{-4} / 4{,}91 \times 10^{-4} = 0{,}51$ m/s
\end{itemize}

Norme : $0{,}5 \leq 0{,}51 \leq 1{,}5$ m/s $\Rightarrow$ \textbf{Oui, la norme est respectée dans les branches} (tout juste).

\smallskip
\textit{Barème : 0,5 pt pour la lecture + 0,5 pt pour la conclusion argumentée.}
\end{corrbox}

% --- Q4 ---
\textbf{Q4} \comp{ANA} — Recherche du diamètre conforme par essais. \pts{1 pt}

\begin{corrbox}
On fait varier $d_1$ pour obtenir $0{,}5 \leq v_1 \leq 1{,}5$ m/s (avec $Q_v = 1{,}8$ m³/h) :

\begin{center}
\renewcommand{\arraystretch}{1.5}
\begin{tabular}{|>{\centering\arraybackslash}m{2.5cm}|>{\centering\arraybackslash}m{2.5cm}|>{\centering\arraybackslash}m{2.5cm}|>{\centering\arraybackslash}m{2.5cm}|}
\hline
\textbf{Essai} & \textbf{Diamètre $d_1$} & \textbf{Vitesse $v_1$} & \textbf{Conforme ?} \\
\hline
1 & 40 mm & \textcolor{rouge}{0,40 m/s} & \textcolor{rouge}{Non (< 0,5)} \\
\hline
2 & \textcolor{rouge}{35 mm} & \textcolor{rouge}{0,52 m/s} & \textcolor{vertpro}{\textbf{Oui}} \\
\hline
3 & \textcolor{rouge}{25 mm} & \textcolor{rouge}{1,02 m/s} & \textcolor{vertpro}{\textbf{Oui}} \\
\hline
\end{tabular}
\renewcommand{\arraystretch}{1}
\end{center}

\textbf{Vérifications :}
\begin{itemize}[nosep]
  \item $d_1 = 35$ mm : $S = \pi \times (0{,}0175)^2 = 9{,}62 \times 10^{-4}$ m² $\Rightarrow v = 5{,}00 \times 10^{-4} / 9{,}62 \times 10^{-4} = 0{,}52$ m/s \checkmark
  \item $d_1 = 25$ mm : $S = \pi \times (0{,}0125)^2 = 4{,}91 \times 10^{-4}$ m² $\Rightarrow v = 5{,}00 \times 10^{-4} / 4{,}91 \times 10^{-4} = 1{,}02$ m/s \checkmark
\end{itemize}

\textbf{Remarque :} à $d_1 = 20$ mm, $v = 1{,}59$ m/s $> 1{,}5$ m/s (trop rapide). La plage conforme est donc entre 22 et 36 mm environ.

\smallskip
\textit{Barème : 0,25 pt par essai correct (valeurs + conclusion).}
\end{corrbox}

% --- Q5 ---
\textbf{Q5} \comp{COM} — Choix et justification du diamètre. \pts{1 pt}

\begin{corrbox}
\textbf{Diamètre choisi :} $d_1 = \boxed{30 \text{ mm}}$

\textbf{Justification :}
\begin{itemize}[nosep]
  \item Avec $d_1 = 30$ mm : $v_1 = 0{,}71$ m/s, bien dans la norme (0,5 – 1,5 m/s)
  \item Un diamètre de 30 mm est plus petit que 35 mm, donc \textbf{moins coûteux} (moins de matière, moins cher à l'achat)
  \item La vitesse de 0,71 m/s offre une bonne marge par rapport aux deux limites (loin de 0,5 et loin de 1,5)
  \item Un diamètre de 25 mm fonctionnerait aussi mais la vitesse (1,02 m/s) est plus élevée, ce qui peut générer davantage de bruit dans les tuyaux
\end{itemize}

Le choix de $d_1 = 30$ mm est un bon \textbf{compromis entre coût et confort acoustique}.

\smallskip
\textit{Barème : 0,5 pt pour le choix d'un diamètre conforme + 0,5 pt pour une justification argumentée (coût + norme).}
\end{corrbox}

\newpage

% ================================================================
\section*{Exercice 2 — Dimensionnement électrique d'une PAC \pts{5 points}}
% ================================================================

\begin{contextebox}
\textbf{Données :} $P = 3\,200$ W \quad $\cos\varphi = 0{,}85$ \quad $U = 230$ V (monophasé, 50 Hz)\\
Fonctionnement : 10 h/jour $\times$ 180 jours. Prix kWh : 0,20 €
\end{contextebox}

\vspace{6pt}

% --- Q1 ---
\textbf{Q1} \comp{REA} — Calculer la puissance apparente $S$. \pts{1 pt}

\begin{corrbox}
On sait que $\cos\varphi = \dfrac{P}{S}$, donc :
\[
S = \frac{P}{\cos\varphi} = \frac{3\,200}{0{,}85} = \boxed{3\,765 \text{ VA}}
\]

\textit{Barème : 0,5 pt pour la formule posée + 0,5 pt pour le résultat avec unité.}
\end{corrbox}

% --- Q2 ---
\textbf{Q2} \comp{REA} — Calculer l'intensité du courant $I$. \pts{1 pt}

\begin{corrbox}
On utilise $S = U \times I$, donc :
\[
I = \frac{S}{U} = \frac{3\,765}{230} = \boxed{16{,}4 \text{ A}}
\]

\textit{Autre méthode directe :}
\[
I = \frac{P}{U \times \cos\varphi} = \frac{3\,200}{230 \times 0{,}85} = \frac{3\,200}{195{,}5} = 16{,}4 \text{ A}
\]

\textit{Barème : 0,5 pt pour la formule + 0,5 pt pour le résultat avec unité.}
\end{corrbox}

% --- Q3 ---
\textbf{Q3} \comp{ANA} — Le disjoncteur de 20 A est-il correctement dimensionné ? \pts{1 pt}

\begin{corrbox}
L'intensité absorbée par la PAC est $I = 16{,}4$ A.

Le disjoncteur est calibré à 20 A.

Or $16{,}4 < 20$ A $\Rightarrow$ \textbf{Oui}, le disjoncteur est correctement dimensionné.

Il reste une marge de $20 - 16{,}4 = 3{,}6$ A, ce qui est suffisant pour les appels de courant au démarrage.

\smallskip
\textit{Barème : 0,5 pt pour la comparaison + 0,5 pt pour la conclusion justifiée.}
\end{corrbox}

% --- Q4 ---
\textbf{Q4} \comp{REA} — Énergie consommée et coût annuel. \pts{1 pt}

\begin{corrbox}
\textbf{Durée de fonctionnement sur la saison :}
\[
t = 10 \times 180 = 1\,800 \text{ h}
\]

\textbf{Énergie consommée :}
\[
E = P \times t = 3\,200 \times 1\,800 = 5\,760\,000 \text{ Wh} = \boxed{5\,760 \text{ kWh}}
\]

\textbf{Coût annuel :}
\[
\text{Coût} = 5\,760 \times 0{,}20 = \boxed{1\,152 \text{ €}}
\]

\textit{Barème : 0,25 pt durée + 0,25 pt énergie en kWh + 0,5 pt coût.}
\end{corrbox}

% --- Q5 ---
\textbf{Q5} \comp{COM} — Expliquer l'impact d'un mauvais facteur de puissance. \pts{1 pt}

\begin{corrbox}
\textbf{Le problème :} la formule $P = U \times I \times \cos\varphi$ montre que, à puissance active $P$ constante :
\[
I = \frac{P}{U \times \cos\varphi}
\]

Si $\cos\varphi$ diminue, l'intensité $I$ \textbf{augmente} pour fournir la même puissance utile.

\textbf{Comparaison chiffrée :}
\begin{itemize}[nosep]
  \item Avec $\cos\varphi = 0{,}85$ : $I = \frac{3\,200}{230 \times 0{,}85} = 16{,}4$ A
  \item Avec $\cos\varphi = 0{,}95$ : $I = \frac{3\,200}{230 \times 0{,}95} = 14{,}6$ A
\end{itemize}

Un mauvais $\cos\varphi$ entraîne un courant plus élevé ($+1{,}8$ A ici), ce qui provoque :
\begin{enumerate}[nosep]
  \item Des \textbf{pertes par effet Joule plus importantes} dans les câbles ($P_J = RI^2$)
  \item Un \textbf{dimensionnement plus coûteux} (câbles, disjoncteurs de calibre supérieur)
  \item Des \textbf{pénalités facturées par le fournisseur} d'électricité pour les professionnels dont le $\cos\varphi$ est trop faible
\end{enumerate}

\textbf{Solution :} installer un condensateur de compensation pour améliorer le $\cos\varphi$ vers 0,95.

\smallskip
\textit{Barème : 0,5 pt pour l'explication (I augmente si $\cos\varphi$ baisse) + 0,5 pt pour les conséquences.}
\end{corrbox}

\vfill

\begin{center}
\textcolor{rouge}{\rule{12cm}{1pt}}\\[8pt]
\begin{tabular}{|c|c|c|c|}
\hline
\textbf{Exercice} & \textbf{Thème} & \textbf{Questions} & \textbf{Points} \\
\hline
1 & Débit de fluide — TP numérique (simulation) & Q1 à Q5 & \textbf{/5} \\
\hline
2 & Puissance en régime sinusoïdal (PAC) & Q1 à Q5 & \textbf{/5} \\
\hline
\multicolumn{3}{|r|}{\textbf{Total}} & \textbf{/10} \\
\hline
\end{tabular}
\end{center}

\end{document}
